\begin{tikzpicture}[
    >=Latex,
    font=\small,
    bloque/.style={
        rectangle,
        draw,
        rounded corners=2pt,
        minimum height=0.9cm,
        minimum width=3.4cm,
        text width=3.2cm,
        align=center,
        inner sep=3pt
    },
    grupo/.style={
        rectangle,
        draw,
        dashed,
        rounded corners=3pt,
        inner sep=7pt
    },
    titulo/.style={
        font=\footnotesize\bfseries,
        anchor=south
    },
    flecha/.style={
        -{Latex[width=6pt,length=5pt]},
        thick
    }
]

% Entradas (x = 0)
\node[bloque] (celda)     at (0, 1.8)  {Celda de carga};
\node[bloque] (hx)        at (0, 0.6)  {Convertidor HX711};
\node[bloque] (teclado)   at (0, -0.6) {Teclado o botones};
\node[bloque] (presencia) at (0, -1.8) {Sensor de presencia del recipiente};

% Procesamiento (x = 5)
\node[bloque, minimum height=3.3cm] (mcu) at (5, 0) {Microcontrolador\\(Arduino o ESP32)};

% Salidas (x = 10 y 14)
\node[bloque] (driver) at (10, 1.2)  {Controlador de motor};
\node[bloque] (alim)   at (14, 1.2)  {Motor del tornillo o del vibrador};
\node[bloque] (servo)  at (10, 0)    {Servomotor de la compuerta};
\node[bloque] (lcd)    at (10, -1.2) {Pantalla LCD};

% Conexiones
\draw[flecha] (celda) -- (hx);
\draw[flecha] (hx.east) -- (hx.east -| mcu.west);
\draw[flecha] (teclado.east) -- (teclado.east -| mcu.west);
\draw[flecha] (presencia.east) -- ++(0.6,0) |- ([yshift=0.35cm]mcu.south west);
\draw[flecha] (mcu.east |- driver.west) -- (driver.west);
\draw[flecha] (driver) -- (alim);
\draw[flecha] (mcu.east |- servo.west) -- (servo.west);
\draw[flecha] (mcu.east |- lcd.west) -- (lcd.west);

% Agrupaciones
\begin{scope}[on background layer]
    \node[grupo, fit=(celda)(presencia)] (gin) {};
    \node[grupo, fit=(mcu)] (gproc) {};
    \node[grupo, fit=(driver)(alim)(lcd)] (gout) {};
\end{scope}
\node[titulo] at (gin.north) {Entradas};
\node[titulo] at (gproc.north) {Procesamiento};
\node[titulo] at (gout.north) {Salidas};

\end{tikzpicture}
